\documentclass[16pt, letterpaper]{article}
\usepackage[T2A]{fontenc}
\usepackage{amsfonts}

\title{42. Теорема единственности для голоморфных функций (с примерами)}
\author{Курбанов Сергей}
\date{Математический анализ, 3 семестр}

\begin{document}

\maketitle

\section*{Теорема}
Пусть $f, g \in \mathcal{A}(D)$, множество $$E = \{z \in D : f(z) = g(z)\}$$ имеет предельную точку в $D$. Тогда $f \equiv g$ в $D$.

\vspace{5mm}

\section*{Доказательство} По условию $\exists$ такая последовательность $\{z_k\}$ точек $D$, что $f(z_k) = g(z_k)$ и $z_k \rightarrow z_0 \in D$. Положим $h = f - g$. Тогда $h \in \mathcal{A}(D)$ и $h(z_k) = 0$ при всех $k \in \mathbb{N}$. по непрерывности $h(z_0) = 0$, т. е. $z_0$ -- не изолированный ноль $h$. По теореме об изолированности нулей голоморфных функций $h \equiv 0$, т. е. $f \equiv g$. $\Box$

\vspace{5mm}

\section*{Примеры}
\textbf{1.} Пусть $D = \mathbb{C}$, $E$ -- невырожденный промежуток в $\mathbb{R}$. Тогда каждая точка $E$ является предельной для $E$ и принадлежит $D$. Поэтому функция $z \mapsto \displaystyle\sum_{k = 0}^\infty \frac{z^k}{k!}$ -- единственная целая функция, совпадающая на $\mathbb{R}$ с экспонентой вещественного аргумента. Аналогично, суммы соответствующих степенных рядов являютсяединственными целыми функциями, продолжающими синус и косинус с промежутка $(0, \frac{\pi}{2})$ на $\mathbb{C}$.

\vspace{5mm}

\textbf{2.} Теорема единственности может использоваться для доказательства тождеств. Например, основное тригонометрическое тождество $$sin^2 z + cos^2 z = 1$$ при $z \in (0, \frac{\pi}{2})$ следует из теоремы Пифагора. По теореме единственности оно верно для всех $z \in \mathbb{C}$, т. к. обе его части -- целые функции. Похожим рассуждением все тождества для показательной и тригонометрических функций распространяются с вещественных значений аргумента на комплексные. В учебнике Виноградова ($\S$ 4 глава 8) они доказывались с помощью операций с рядами и формул Эйлера.

\vspace{5mm}

\textbf{3.} Условие принадлежности предельной точки множеству $D$ в теореме единственности существенно. Так, целая функция $sin$ отлична от тождественного нуля и обращается в ноль на счетном множестве $\{k\pi : k \in \mathbb{Z}\}$, которое не имеет предельной точки в $\mathbb{C}$. Функция $z \mapsto sin \frac{1}{z}$ голоморфна в круге $B(1, 1)$, отлична от тождественного нуля и обращается в нольна счетном множестве $E = \{\frac{1}{k\pi} : k \in \mathbb{N}\}$, единственная предельная точка которого $0$ не принадлежит $B(1, 1)$. Функции же, голоморфной в круге $B(0, 1)$, обращающейся в ноль на множестве $E$ и отличной от тождественного нуля, по теореме единственности не $\exists$.

\end{document}

